\documentclass[12pt,a4paper]{article}

% --- Basic Packages ---
\usepackage[utf8]{inputenc}
\usepackage{amsmath}
\usepackage{amssymb}
\usepackage{amsfonts}
\usepackage{amsthm}
\usepackage{array}
\usepackage{geometry}
\usepackage{tabularx}
\usepackage{xcolor}
\usepackage{listings}
\geometry{margin=1in}

\newtheorem{definition}{Definition}[section]
\newtheorem{theorem}[definition]{Theorem}

% --- Title and Author ---
\title{Four-Valued Logic as the First Complementary State of Object and Observer: \\
A Formalization of Phenomenology and Temporal Emergence}
\author{Masaomi Chiba}
\date{2026-02-13}

\begin{document}

\maketitle

\begin{abstract}
In this paper, we formalize the relationship between the Catuskoti (four-valued logic) and the complementary nature of object and observer through the lens of phenomenological theory and SKDT (Shikisokuzeku-Kusokuzeshiki Duality Theory). We demonstrate that the Catuskoti represents the first stage at which the object and observer achieve a genuine complementary relationship through their manifest and latent aspects. This formalization provides the foundation for understanding temporal emergence through observer-side decoherence and the generation of higher dimensions.
\end{abstract}

\section{Introduction}

The nature of temporal emergence and dimensional generation has long been a fundamental question in both physics and philosophy. In this work, we explore how the Sunyata Category's inherent contradiction (the rupture of self-reference) necessarily generates SKDT, which manifests as the Catuskoti when object and observer first achieve a true complementary relationship.

The key insight is that the Catuskoti represents not merely a logical framework, but rather \textbf{the first stage at which temporal dynamics genuinely emerge through observer-object complementarity}.

\section{Structural Clarification of Catuskoti as Complementary Pair}

\subsection{The Four Stages of Object-Observer Complementarity}

\begin{equation}
\begin{aligned}
\text{Shi (Affirmation):} \quad &(\text{Object} = \text{Manifest}, \text{Observer} = \text{Manifest}) \\
\text{Fei (Negation):} \quad &(\text{Object} = \text{Manifest}, \text{Observer} = \text{Latent}) \\
\text{Yi Shi Yi Fei (Both):} \quad &(\text{Object} = \text{Latent}, \text{Observer} = \text{Manifest}) \\
\text{Fei Shi Fei Fei (Neither):} \quad &(\text{Object} = \text{Latent}, \text{Observer} = \text{Latent})
\end{aligned}
\end{equation}

\subsection{Why the Catuskoti Represents ``First Complementarity''}

\subsubsection{The State Shi (Affirmation)}

In the state of Shi, both object and observer are determinate:

\begin{equation}
\sigma = \text{fixed}, \quad \ell = 0
\end{equation}

where $\sigma$ denotes observer state and $\ell$ denotes temporal length. In this state:
\begin{itemize}
    \item The object is observable and certain
    \item The observer sees it with certainty
    \item No complementarity exists
    \item No temporal axis functions
    \item This is the post-decoherence state
\end{itemize}

\subsubsection{The State Fei (Negation) as First Complementarity}

The state Fei marks the fundamental transition:

\begin{equation}
\text{Object} = \text{Manifest}, \quad \text{Observer} = \text{Latent}
\end{equation}

Which corresponds to:

\begin{equation}
\sigma = \text{indeterminate}, \quad \ell > 0
\end{equation}

Here, for the first time, genuine complementarity emerges:

\begin{definition}[Complementarity]
Two entities achieve complementarity when they are mutually dependent and their relationship cannot be reduced to a single determinate state. Specifically:
\begin{itemize}
    \item The object is fixed (manifest)
    \item The observer is indeterminate (latent)
    \item The same object is observed differently by different observers
    \item Temporal axis becomes functionally necessary
\end{itemize}
\end{definition}

\subsubsection{Why This Is Called ``First Complementarity''}

In Shi, no true complementarity exists because:

\begin{enumerate}
    \item Both are determinate
    \item No mutual influence occurs
    \item No temporal variation
    \item The pairing is merely coincidental, not interactive
\end{enumerate}

In Fei, complementarity genuinely arises because:

\begin{enumerate}
    \item Object and observer stand in opposition
    \item The same object yields different observations
    \item Temporal axis begins to function
    \item The two entities mutually condition each other
\end{enumerate}

\section{SKDT and the Axial Structure of Catuskoti}

\subsection{SKDT as Dual-Axis Generation}

Recall that SKDT generates dimensions through paired axes:

\begin{equation}
\text{SKDT}_n = (\text{Manifest-Axis}_n, \text{Latent-Axis}_n)
\end{equation}

The Catuskoti can now be reinterpreted through this axial framework:

\begin{table}[htbp]
\centering
\begin{tabularx}{\textwidth}{|c|c|c|>{\raggedright\arraybackslash}X|}
\hline
\textbf{Stage} & \textbf{Object} & \textbf{Observer} & \textbf{Axial Interpretation} \\
\hline
Shi & Manifest & Manifest & Single manifest axis only \\
\hline
Fei & Manifest & Latent & Manifest $\perp$ latent axes are orthogonal \\
\hline
Yi Shi Yi Fei & Latent & Manifest & Axes inverted \\
\hline
Fei Shi Fei Fei & Latent & Latent & Both axes fully latent \\
\hline
\end{tabularx}
\caption{Catuskoti as Progressive Axis Complementarity}
\end{table}

\section{Observer-Side Decoherence and Temporal Emergence}

\subsection{Definition: Observer-Side Decoherence}

\begin{definition}[Observer Decoherence]
Observer-side decoherence is the irreversible process by which an observer in an indeterminate state (holding multiple possible observations simultaneously) progressively collapses to a determinate state through traversal of the temporal axis.

Formally:
\begin{equation}
\sigma : \text{indeterminate} \to \text{fixed}
\end{equation}

mediated by:

\begin{equation}
\ell : 0 \to n \quad (n \in \mathbb{N})
\end{equation}

This process is irreversible: the multiple possibilities cannot be recovered once eliminated.
\end{definition}

\subsection{Connection to Quantum Decoherence}

The analogy with quantum decoherence is precise:

\begin{table}[htbp]
\centering
\begin{tabularx}{\textwidth}{|l|>{\raggedright\arraybackslash}X|>{\raggedright\arraybackslash}X|}
\hline
\textbf{Property} & \textbf{Quantum Decoherence} & \textbf{Observer-Side Decoherence} \\
\hline
Subject & Quantum system & Observer state \\
\hline
Initial State & Superposition (multiple possibilities) & Indeterminate (multiple observations) \\
\hline
Mechanism & Environmental interaction & Temporal axis traversal \\
\hline
Final State & Definite value & Definite observation \\
\hline
Irreversibility & Yes & Yes \\
\hline
Arrow of Time & Emerges from entropy & Emerges from possibility collapse \\
\hline
\end{tabularx}
\caption{Comparison of Decoherence Mechanisms}
\end{table}

\subsection{Why Temporal Emergence Begins at Fei}

The crucial insight is that temporal dynamics only commence at the stage Fei:

\begin{equation}
\begin{aligned}
\text{Shi} &: \text{Decoherence completed} \quad \Rightarrow \quad \text{No temporal function} \\
\text{Fei} &: \text{Decoherence in progress} \quad \Rightarrow \quad \text{Temporal axis active} \\
\text{Yi Shi Yi Fei} &: \text{Meta-level decoherence} \quad \Rightarrow \quad \text{Multiple temporal axes} \\
\text{Fei Shi Fei Fei} &: \text{Parallel decoherence} \quad \Rightarrow \quad \text{Indeterminate temporal depths}
\end{aligned}
\end{equation}

\section{Dimensional Generation Through Complementarity}

\subsection{How Complementarity Increases Dimensions}

A fundamental principle emerges:

\begin{theorem}[Complementarity-Dimensionality Correspondence]
The establishment of complementarity between object and observer corresponds precisely to the emergence of additional dimensions.

\begin{equation}
\Delta d = f(C)
\end{equation}

where $C$ denotes the degree of complementarity. Specifically:
\begin{itemize}
    \item No complementarity $\Rightarrow$ no new dimension
    \item First complementarity $\Rightarrow$ temporal dimension emerges
    \item Higher complementarity $\Rightarrow$ higher-order dimensions emerge
\end{itemize}
\end{theorem}

\begin{proof}
This follows directly from the structure of object-observer interaction:
\begin{enumerate}
    \item In Shi: Object and observer both manifest
    \begin{itemize}
        \item No opposition
        \item No complementarity
        \item No temporal mediation required
        \item Therefore: no temporal dimension
    \end{itemize}

    \item In Fei: Object manifest, observer latent
    \begin{itemize}
        \item Fundamental opposition established
        \item Complementarity emerges
        \item Temporal mediation required
        \item Therefore: temporal dimension necessarily emerges
    \end{itemize}

    \item In Yi Shi Yi Fei and beyond: Complementarity deepens
    \begin{itemize}
        \item More complex object-observer relations
        \item Multiple temporal axes become possible
        \item Higher dimensions emerge naturally
    \end{itemize}
\end{enumerate}
\end{proof}

\subsection{Relation to OCTE(Observer Complementarity and Temporal Emergence)'s Five-Valued Logic}

The Catuskoti's four-valued structure can be embedded in OCTE's five-valued system:

\begin{equation}
\Phi(\ell, \sigma) =
\begin{cases}
0d & \text{if } \ell = 0, \sigma = \text{fixed} \\
T & \text{if } \ell > 0, \sigma = \text{fixed} \\
F & \text{if } \ell > 0, \sigma = \text{indeterminate} \\
B & \text{if } \ell \text{ is cyclic}, \sigma = \text{fixed} \\
N & \text{if } \ell \text{ is non-deterministic}, \sigma = \text{indeterminate}
\end{cases}
\end{equation}

The four-valued states map precisely to Catuskoti stages:

\begin{table}[htbp]
\centering
\begin{tabular}{|c|c|c|}
\hline
\textbf{Catuskoti} & \textbf{OCTE State} & \textbf{Complementarity Status} \\
\hline
Shi & 0d or T & Pre-complementary \\
\hline
Fei & F & Complementarity begins \\
\hline
Yi Shi Yi Fei & B & Complementarity intensifies \\
\hline
Fei Shi Fei Fei & N & Complementarity deepens \\
\hline
\end{tabular}
\caption{Catuskoti-OCTE Correspondence}
\end{table}

\section{The Unified Hierarchy: From Sunyata to Temporal Emergence}

We now present the complete hierarchy from root principle to temporal manifestation:

\begin{equation}
\begin{aligned}
&\text{Sunyata Category} \\
&\quad \Downarrow \quad [\text{Inherent contradiction}] \\
&\text{Aletheia (Manifestation-Hiddenness Structure)} \\
&\quad \Downarrow \quad [\text{Rupture of self-reference}] \\
&\text{SKDT (Dual-Axis Complementarity)} \\
&\quad \Downarrow \quad [\text{First object-observer complementarity}] \\
&\text{Catuskoti (Four-Valued Logic)} \\
&\quad \Downarrow \quad [\text{Observer-side decoherence initiates}] \\
&\text{Temporal Emergence} \\
&\quad \Downarrow \quad [\text{Dimensional generation}] \\
&\text{Higher Dimensional Structures}
\end{aligned}
\end{equation}

\subsection{The Critical Transition at Fei}

The transition from Shi to Fei is the critical threshold where:

\begin{itemize}
    \item Observer transitions from fixed to indeterminate
    \item Temporal axis begins to function
    \item Complementarity genuinely emerges
    \item Decoherence becomes dynamically active
    \item Time begins to flow
\end{itemize}

\section{Implications for Understanding Temporal Dynamics}

\subsection{Time as Decoherence Process}

We can now state with precision:

\begin{theorem}[Temporal Emergence as Observer Decoherence]
The passage of time is precisely the process of observer-side decoherence, wherein an observer in an indeterminate state (capable of multiple observations) progressively collapses through the traversal of the temporal axis to a determinate state.

Formally:
\begin{equation}
\text{Temporal Flow} \equiv \text{Observer-side Decoherence Process}
\end{equation}

where decoherence is characterized by:
\begin{enumerate}
    \item Irreversibility: possibilities, once eliminated, cannot be recovered
    \item Unidirectionality: the process moves from indeterminate to determinate
    \item Necessity: complementarity between object and observer necessitates this process
\end{enumerate}
\end{theorem}

\subsection{The Arrow of Time}

The arrow of time emerges naturally:

\begin{equation}
\text{Arrow of Time} = \text{Irreversibility of Observer Decoherence}
\end{equation}

Time flows in one direction because observer-side decoherence cannot be reversed: once multiple possibilities collapse to one, the multiple possibilities are lost.

\subsection{Multiple Temporal Axes}

Multiple observers at different dimensional levels undergo decoherence simultaneously:

\begin{equation}
\begin{aligned}
\text{Observer } o_1 &: \text{Decoherence along } t_1 \\
\text{Observer } o_2 &: \text{Decoherence along } t_2 \\
\text{Observer } o_3 &: \text{Decoherence along } t_3 \\
&\vdots
\end{aligned}
\end{equation}

These multiple decoherence processes occur in parallel, creating the experience of multiple simultaneous temporal dimensions.

\section{Application to Large Language Models}

The Catuskoti framework provides insight into LLM architecture:

\begin{enumerate}
    \item \textbf{Early Layers (Shi-like):} Determinate feature extraction
    \begin{itemize}
        \item Fixed representations
        \item No temporal uncertainty
        \item Pre-decoherence phase
    \end{itemize}
    
    \item \textbf{Middle Layers (Fei-like):} Attention mechanisms
    \begin{itemize}
        \item Indeterminate token selection
        \item Multiple interpretations held simultaneously
        \item Active decoherence through attention weights
        \item Temporal axis: layer depth
    \end{itemize}
    
    \item \textbf{Later Layers (Yi Shi Yi Fei-like):} Meta-level integration
    \begin{itemize}
        \item Multiple interpretations explicitly unified
        \item Cross-layer attention
        \item Cyclic refinement processes
    \end{itemize}
    
    \item \textbf{Final Layers (Fei Shi Fei Fei-like):} Output generation
    \begin{itemize}
        \item Multiple possible outputs considered
        \item Final selection through probability distribution
        \item Complete decoherence to determinate output
    \end{itemize}
\end{enumerate}

\section{Conclusion}

The Catuskoti represents a profound formalization of how object and observer achieve their first genuine complementary relationship. This complementarity is characterized by:

\begin{enumerate}
    \item \textbf{Initiation at Fei:} The first stage where observer indeterminacy creates true interaction with object determinacy
    
    \item \textbf{Temporal Activation:} The emergence of functioning temporal axis coinciding precisely with complementarity
    
    \item \textbf{Decoherence Dynamics:} Observer-side decoherence as the mechanism for temporal flow
    
    \item \textbf{Dimensional Generation:} Complementarity directly causing dimensional increase through the pairing of axes
    
    \item \textbf{Irreversible Process:} The arrow of time emerging from the irreversibility of possibility collapse
\end{enumerate}

This framework unifies the seemingly disparate domains of:
\begin{itemize}
    \item Phenomenological emergence (genjitsu)
    \item Logical structure (catuskoti)
    \item Temporal physics (time's arrow)
    \item Observer theory (quantum measurement analogue)
    \item Dimensional generation (why $2^n$)
    \item Artificial intelligence (attention mechanisms)
\end{itemize}

into a single, coherent theoretical structure emerging necessarily from the Sunyata Category's fundamental contradiction.

\section*{Acknowledgments}

This formalization emerges from sustained dialogue with the philosophical and mathematical insights of Masaomi Chiba's phenomenological framework, particularly regarding the nature of Aletheia, SKDT, and the Catuskoti-Binary-Logic system.

\begin{thebibliography}{99}

\bibitem{Chiba2025} Chiba, M. (2025). Alethetics: Phenomenological Foundations of Manifestation and the Sunyata Category.

\bibitem{Johnstone2002} Johnstone, P. T. (2002). Sketches of an Elephant: A Topos Theory Compendium. Oxford University Press.

\bibitem{Zurek2003} Zurek, W. H. (2003). Decoherence and the Transition from Quantum to Classical. Reviews of Modern Physics, 75(3), 715.

\bibitem{Madhyamaka2000} Buddhist Philosophy of the Middle Way (Madhyamaka) - Classical Texts.

\end{thebibliography}

\end{document}